\begin{table}[H]
\centering
\caption{Productividad y remuneración por departamento, 2024pr}
\label{tab:dept_productividad_remuneracion_24}
\scriptsize
\begin{tabular}{lrrrrr}
\toprule
Departamento & PIB/hora 2024pr & Rem./hora 2024 & Residuo & PIB/trab. 2024pr & Rem./trab. 2024 \\
\midrule
Bogotá D.C. & 28,6 & 15,0 & 30,4\% & 67,4 & 2,95 \\
Caldas & 15,6 & 9,4 & 21,4\% & 36,6 & 1,85 \\
Quindío & 15,1 & 9,1 & 19,5\% & 33,8 & 1,71 \\
Norte de Santander & 10,0 & 7,4 & 18,8\% & 24,3 & 1,49 \\
Antioquia & 20,3 & 10,9 & 18,7\% & 47,7 & 2,12 \\
Caquetá & 10,5 & 7,4 & 17,6\% & 24,9 & 1,47 \\
Risaralda & 16,4 & 9,2 & 15,3\% & 38,7 & 1,82 \\
Cundinamarca & 17,4 & 9,2 & 11,2\% & 40,4 & 1,79 \\
Nariño & 9,1 & 6,5 & 9,4\% & 17,8 & 1,06 \\
Atlántico & 16,0 & 8,1 & 2,8\% & 37,3 & 1,58 \\
Huila & 14,2 & 7,3 & -0,8\% & 31,6 & 1,36 \\
Córdoba & 10,7 & 6,1 & -4,8\% & 22,0 & 1,04 \\
Magdalena & 10,3 & 5,9 & -5,1\% & 23,9 & 1,15 \\
Tolima & 16,9 & 7,6 & -7,2\% & 40,0 & 1,50 \\
Valle del Cauca & 23,4 & 9,3 & -7,7\% & 54,5 & 1,79 \\
Sucre & 10,4 & 5,7 & -9,3\% & 22,0 & 1,01 \\
Cauca & 13,7 & 6,5 & -9,8\% & 26,6 & 1,06 \\
Cesar & 13,8 & 6,4 & -11,8\% & 32,6 & 1,26 \\
Boyacá & 22,6 & 8,5 & -13,3\% & 47,5 & 1,49 \\
Chocó & 11,9 & 5,8 & -14,0\% & 28,6 & 1,16 \\
Santander & 26,1 & 9,1 & -15,8\% & 62,0 & 1,80 \\
Meta & 25,9 & 9,0 & -15,9\% & 63,1 & 1,83 \\
Bolívar & 19,8 & 7,0 & -22,1\% & 45,1 & 1,33 \\
La Guajira & 11,1 & 4,8 & -26,4\% & 24,8 & 0,89 \\
\bottomrule
\end{tabular}
\caption*{\footnotesize Nota: PIB por hora en miles de pesos constantes de 2015; remuneración por hora en miles de pesos constantes de 2025; PIB por trabajador en millones de pesos constantes de 2015; remuneración por trabajador en millones de pesos constantes de 2025 al mes. El residuo mide la distancia porcentual frente a la tendencia lineal entre PIB por hora y remuneración por hora. Fuente: cálculos propios con DANE y GEIH.}
\end{table}